\section{AI 技术方案}

\subsection{AI 能力全景}

系统没有把“AI”做成一个孤立的聊天框，而是把它拆成六种职责明确的能力。每种能力都有独立的模型绑定、输入输出契约与失败降级路径；任何一项不可用，都不会让整个产品变成空白页。

\begin{table}[htbp]
  \centering
  \small
  \caption{AI 能力、模型角色与降级路径}
  \label{tab:ai-map}
  \begin{tabularx}{\textwidth}{L{2.3cm}L{2.1cm}L{4.3cm}L{2.7cm}Y}
    \toprule
\textbf{能力} & \textbf{模型角色} & \textbf{输入 → 输出} & \textbf{失败时} & \textbf{说明} \\
    \midrule
    关系判定 & \code{relation} & 论文对 + rubric + 反例 → 多标签 JSON & 启发式打分，界面标注“启发式” & 只在硬约束过滤后的候选上调用 \\
    文献问答 & \code{ask} & 证据块 + 问题 → 流式文本 & 保留检索摘录；未绑定模型则明确提示 & 资料事实带角标，通用知识不加角标 \\
    查询规划 & \code{ask} & 星空摘要 + 问题 → 工具步骤 JSON & 退回默认混合检索 & 最多 3 步，参数再过 schema \\
    向量召回 & \code{embedding} & 文本块 / 查询 → 单位向量 & 退回 BM25 词面检索 & 失败不是错误，诊断中标注 \code{lexical} \\
    图像理解 & \code{ask}（视觉） & 图片 → OCR/表格/图表/数值/区域 & 不使用图像证据，仅文本问答 & 最多 4 张，坐标归一化到 0--1 \\
    标题翻译 & \code{translate} & 英文标题/摘要 → 中文 & 保留原文 & 独立模型，避免占用问答额度 \\
    \bottomrule
  \end{tabularx}
\end{table}

\subsection{选型理由：为什么是“通用模型 + 结构化编排”}

\subsubsection{不微调：关系分类的边界是用户可编辑的}

六类关系不是固定标签，而是可编辑的判定标准：用户可以修改“基线”的引用条件、年份差、产出上限，也可以新建最多 24 条自定义关系。若把关系分类微调进模型，每次改 rubric 都要重新标注与训练，且模型无法解释“为什么这条边成立”。本项目选择“通用大模型 + 可编辑 rubric + 反例 + 原文证据”的零样本方案：改一条规则只需要改一个 JSON 对象，判定结果仍能给出理由与引用片段。代价是单次判定需要付提示词成本，这一点通过硬约束预过滤、批量判定与缓存来抵消。

\subsubsection{不绑定 function calling：JSON 模式是更通用的能力}

不同供应商对 function calling 的支持参差不齐，流式工具调用的增量格式也不一致。项目把“选工具”收敛成一次结构化 JSON 输出，再由服务端用 Zod 校验、用白名单过滤。这样即使供应商只支持最基础的 chat completions，规划器仍然可用；解析失败时也能整体退回默认检索，不会出现“半个工具调用挂在半空”的状态。

\subsubsection{不上向量数据库：语料规模与运维成本不匹配}

每片星空的语料是“几十到数百篇论文的摘要 + 关系证据”，规模在数千个文本块以内。项目在进程内构建 BM25 索引与向量矩阵，按星空版本缓存；同一片星空重复提问不重复付费。引入独立向量数据库会增加部署组件、网络往返与一致性维护，而当前规模下收益有限。代码把召回封装成 \code{EmbeddingProvider} 与 \code{retrieveSourcesDetailed}，未来语料上量时可以替换为外部向量库，而不改上层 Agent。

\subsubsection{不自研 Agent 框架：闭环要可审计}

项目没有引入 LangChain 之类的通用 Agent 框架，而是实现了一个很小的闭环：规划 → 执行 → 证据覆盖评估 → 补查或停止。原因是赛题要求“每个决策都能解释”，而通用框架的中间状态往往藏在回调里。自研闭环把每一步写成 \code{AskStep} 与 \code{AskDecision}，停止条件由确定性规则给出，既能在界面展示轨迹，也能在 \code{pnpm check:agent} 中直接测试生产代码路径。

\subsubsection{不让一个模型包打天下：按任务拆分模型}

关系判定需要稳定、低温、严格 JSON；问答需要更强的语言组织与长上下文；向量模型必须与文本模型解耦；翻译模型不需要很贵。系统按服务分别绑定模型，用户可以在设置页为“判定 / 问答 / 翻译 / 向量”选择不同厂商与型号。这样既能让便宜模型承担高频判定，也能让强模型只在最终回答时出场。

\begin{table}[htbp]
  \centering
  \small
  \caption{关键选型的备选方案对比}
  \label{tab:ai-choice}
  \begin{tabularx}{\textwidth}{L{3.0cm}L{4.6cm}L{4.2cm}Y}
    \toprule
\textbf{决策} & \textbf{备选方案} & \textbf{主要问题} & \textbf{本项目结论} \\
    \midrule
    关系分类 & 微调专用分类器 & 需要标注数据；rubric 变更即失效；难解释 & 通用模型 + rubric + 反例 + 证据核验 \\
    工具调用 & 供应商 function calling & 支持不一；流式增量格式不统一 & 一次性 JSON 规划 + 服务端 schema 校验 \\
    向量检索 & 独立向量数据库 & 多一个运维组件；当前语料规模收益有限 & 进程内索引 + 版本缓存，接口预留替换 \\
    Agent & 通用 Agent 框架 & 中间状态不可控；调试与复现成本高 & 五工具、两轮、五步的自研闭环 \\
    模型 & 单一模型全包 & 成本高；任务间相互牵制 & 判定/问答/翻译/向量分角色绑定 \\
    \bottomrule
  \end{tabularx}
\end{table}

\subsection{集成方式：按账号绑定、服务端解密、统一客户端}

模型服务不读取部署级环境变量。每个账号登录后，在“设置 → 模型服务”里添加自己的服务商、API Key 与模型；服务端用 AES-256-GCM 加密后落库，对外只返回 \code{hasApiKey} 与尾号。调用时由 \code{resolveServiceModel(ownerId, service)} 解析出当前账号的目标模型，密钥只在服务端内存中短时存在。这条设计同时满足三件事：多厂商可切换、密钥不进入浏览器包、不同账号的额度与账单互不污染。

大模型客户端统一封装在 \code{src/server/llm/client.ts}：兼容 OpenAI 兼容接口与 Anthropic Messages API；\code{chatJson} 负责非流式结构化输出，\code{streamChat} 把两种供应商的 SSE 事件收敛成同一种 delta；所有输出都当作不可信输入，JSON 必须过 Zod。客户端还负责超时、重试、错误翻译与用量记账：默认超时 60 秒，向量请求 20 秒，重试 1 次；401/402/403/404/429 会被翻译成“密钥无效 / 余额不足 / 无权限 / 模型不存在 / 限流”的可读提示，且绝不回显请求头。

问答链路端到端使用 SSE。浏览器发起 \code{POST /api/ask}，服务端先做限流、预算、输入护栏与模型解析，再创建 \code{ReadableStream}，把规划轨迹、检索诊断、回答增量与最终引用逐条推给前端。客户端关闭页面或点“停止”时，\code{req.signal} 会中止上游请求，避免继续烧 token。反向代理必须关闭缓冲（Caddy 的 \code{flush_interval -1} 或 Nginx 的 \code{proxy_buffering off}），否则回答会变成最后一次性出现。

\begin{designbox}{为什么把“模型选择”放在账号而不是环境变量}
部署级密钥看起来省事，但会带来三个问题：所有用户共享同一份额度，无法按账号计量；演示环境一旦泄露，影响所有使用者；赛题要求“多厂商可切换”，环境变量只能绑定一个。按账号绑定后，每个用户对自己的密钥与账单负责，系统只保存密文与尾号，既符合自托管的现实，也让“模型服务”成为可演示的产品能力。
\end{designbox}

\subsection{关系判定：两段式管道与多标签 rubric}

关系判定是整条 AI 链路里最容易产生幻觉的环节，也是成本最高的环节。系统把它拆成“确定性候选生成 + 硬约束过滤 + 多标签判定 + 证据核验”四段，顺序不能反。

\begin{figure}[htbp]
\centering
\begin{tikzpicture}[font=\footnotesize,>=Latex,
  stage/.style={draw=rule,fill=white,rounded corners=1.4mm,align=center,inner sep=3pt,minimum height=9mm,text width=2.35cm},
  side/.style={font=\scriptsize,color=muted,align=left,text width=3.1cm},
  arr/.style={->,thick,accent!65}]
  \node[stage] (a) at (0,0) {候选生成\\引用 + 向量 + 概念};
  \node[stage] (b) at (3.0,0) {硬约束过滤\\年份/互引/任务/相似度};
  \node[stage] (c) at (6.0,0) {多标签判定\\一对论文一次调用};
  \node[stage] (d) at (9.0,0) {批处理与缓存\\4 对/批，promptVersion};
  \node[stage] (e) at (12.0,0) {校验与入库\\白名单/证据/不变式};
  \foreach \x/\y in {a/b,b/c,c/d,d/e} \draw[arr] (\x) -- (\y);
  \node[side,anchor=north west] at (-1.15,-1.0) {候选量级从 $O(n^2)$ 降到 $O(n\cdot k)$；这一步不花 token，也不会幻觉。};
  \node[side,anchor=north west] at (8.0,-1.0) {解析失败退回逐对多标签；没有模型时降级为启发式，并在界面标注。};
\end{tikzpicture}
\caption{关系判定管道：确定性过滤在前，模型判定在后，证据核验在最后。}
\label{fig:relation-pipe}
\end{figure}

\subsubsection{候选生成与硬约束}

候选来自三类信号：引文邻域（引用与被引用）、向量近邻（摘要语义相近）、概念共现（共享主题/任务）。硬约束完全确定性：年份差、互引检查、共享任务、文本相似度下限。它把候选量级从 $O(n^2)$ 砍到可接受的 $O(n\cdot k)$，而且这一步是唯一能保证“不出现明显错误”的步骤。不同关系有不同约束：基线要求存在引用且年份差不超过 8 年；同期要求互不引用且年份差不超过 2 年；思路类似要求标题+摘要 TF-IDF 余弦不低于 0.15；同任务异方法要求共享任务。

阈值不是拍脑袋定的。思路类似的阈值最初定在 0.35，实测把整类关系筛成了空集——两篇摘要即使讲同一件事，因用词与句式不同，余弦通常也只在 0.2--0.4 之间。阈值必须贴着指标的尺度定，而不是照着“感觉上应该像”的直觉定。

\subsubsection{多标签提示词：一次比较所有可判关系}

同一对论文可能同时是“基线”与“思路类似”。如果每类关系各问一次，会把两段摘要和大部分 rubric 重复送好几遍。系统按“论文对”送审：一次调用带上这一对可判关系的全部 rubric、反例与结构化事实，返回所有成立的关系。提示词骨架如下。

\begin{lstlisting}
SYSTEM
You are a relation judge for academic papers.
For each pair, decide which of the allowed relations hold.
A relation is valid only if the quoted evidence appears in the paper text.
Never infer a relation from general knowledge alone.

RUBRIC (per relation)
  relationId: baseline | contemporary | similar | alternative | extends | contradicts
  definition: <user-editable definition>
  negativeExamples: <what does NOT count>
  hardFacts: { citesIJ, citesJI, yearGap, sharedTopics, textSimilarity }

INPUT
  paperA: { title, year, abstract (<=400 chars) }
  paperB: { title, year, abstract (<=400 chars) }

OUTPUT (JSON only)
{"results":[{"index":0,"labels":[
  {"relationId":"baseline","confidence":0.82,
   "rationale":"B explicitly compares against A.",
   "evidence":[{"paperId":"A","quote":"We compare against ..."}]}
]}]}
\end{lstlisting}

反例是提示词里最重要的部分。只给正例，模型不知道边界在哪，就会把所有“看起来差不多”的论文对判成同一种关系；“同期”与“思路类似”恰恰最容易互相污染。项目为每条关系写明“什么不算”：只是被顺带引用不算基线；共享一个通用组件不算思路类似；同一方法的两个微调版本不算同任务异方法。这些反例与正例一样是用户可编辑的数据，而不是写死在代码分支里。

\subsubsection{批处理、缓存与成本}

\begin{table}[htbp]
  \centering
  \small
  \caption{关系判定的关键参数与取值理由}
  \label{tab:relation-params}
  \begin{tabularx}{\textwidth}{L{3.3cm}C{1.6cm}Y}
    \toprule
\textbf{参数} & \textbf{取值} & \textbf{理由} \\
    \midrule
    每批论文对 & 4 & 固定 rubric 开销摊薄到 1/4；再大容易输出截断或漏项，整批重来的代价更高 \\
    并发判定数 & 4 & 实测既能压住墙钟时间，又不容易触发供应商限流 \\
    摘要长度上限 & 400 字符 & 摘要第一段通常是“问题 + 方法 + 结果”；从 700 压到 400，输入 token 明显下降，判定质量无可感知变化 \\
    邻域候选池 & 40 & 池子要足够宽才能找到跨分支工作；抓取免费且按种子缓存，判定才花钱 \\
    每类关系送审对数 & 16 & 成本与覆盖的折中；调大能长出更密的语义图，代价接近线性 \\
    传递扩展层数 & 2 & “图能长大”与“成本可控”的折中 \\
    每条关系产出上限 & 6--8 & 防止单一关系淹没画面；超出部分按置信度截断 \\
    \bottomrule
  \end{tabularx}
\end{table}

成本可以粗略写成 $\text{calls} \approx \lceil P \times D / B \rceil$，其中 $P$ 为启用关系数，$D$ 为每类送审对数，$B$ 为批大小。以 6 条关系、16 对、4 对/批估算，一次冷启动的量级约为 24 次调用；实际会因候选不足、缓存命中与“没有成立关系”的空结论缓存而下降。缓存键是“无序论文对 + 关系 id + 该关系的 \code{promptVersion}”：用户只改“基线”的 rubric 时，同一对论文上“扩展”的旧结论仍然可用，不会被误作废；没有成立关系的空结论也会缓存，避免重复送审。

\subsubsection{校验与入库}

模型返回后，服务端按 \code{relationId} 白名单归一化，核验每条引用片段是否能在对应论文文本中找到；找不到的片段计入 \code{rejectedQuotes}，这条关系被丢弃。随后检查不变式：语义边至少一条证据；同一对节点同一关系只保留置信度最高的一条并合并证据；聚合节点不得作为引用端点。违反即拒绝入库，而不是“先存下来再说”。关系报告会记录候选数、判定数、命中数、调用数、估算 token、缓存命中、批处理回退与 \code{rejectedQuotes}——这个被拒片段数本身就是反幻觉能力的证据。

没有大模型密钥时，第二段降级为启发式打分：引用结构决定基线与扩展，主题重叠决定同任务异方法，文本相似度决定思路类似。界面会明确显示“启发式判定”，而不是假装有 AI 结论。这是产品诚信的一部分，也是演示时最容易被追问的地方。

\subsection{RAG：混合召回、图扩展与意图重排}

RAG 的目标不是“把摘要塞进提示词”，而是让回答里的每个事实都有落点。系统先构建语料：每篇论文的标题、摘要、关系证据被切成文本块，每块携带 \code{paperId}、\code{section} 与 \code{locator}；图像证据则以单独的块进入同一套检索与覆盖判定。

\begin{table}[htbp]
  \centering
  \small
  \caption{四类漏召回问题与对应策略}
  \label{tab:rag-problems}
  \begin{tabularx}{\textwidth}{L{4.2cm}Y}
    \toprule
\textbf{问题} & \textbf{策略} \\
    \midrule
    中英混问 / 同义表达 & 查询规划补英文同义词；词面与向量同时召回，避免只靠一种信号 \\
    比较三篇以上 & 先按论文拆子问题，再按多锚点补关系证据；避免第一篇论文占满名额 \\
    两路分数不可比 & BM25 分数无上界、余弦在 $[-1,1]$，直接相加会被 BM25 支配；先变名次，再用 RRF 融合 \\
    召回多但顺序差 & 最后按问题覆盖、论文命中、章节意图与图证据重排 \\
    \bottomrule
  \end{tabularx}
\end{table}

RRF 只看名次，不看分数：$\mathrm{RRF}(d)=\sum_i w_i/(60+\mathrm{rank}_i(d))$。两路召回各自给出排名，融合后再做图结构扩展：如果问题提到多篇论文，系统会从关系边中补入连接这些论文的证据块；扩展块的分数门槛是当前最高分的 0.55 倍，避免弱相关边挤掉正文。重排分数由五项组成：归一化检索分 0.44、问题概念覆盖 0.26、标题命中 0.10、提到论文命中 0.10、章节意图加成；提到论文但证据未覆盖时扣 0.12，图扩展块加 0.04。

置信度与停止门槛是确定性的：$\text{confidence}=0.65\times\text{coverage}+0.35\times\text{normalizedTop}$；当且仅当“至少 2 个来源、覆盖率不低于 0.28、问题提到的论文全部被覆盖、置信度不低于 0.38”时，默认检索被判定为“足够”。这套门槛不是为了让模型“感觉良好”，而是为了让 Agent 知道什么时候不该继续花钱补查。向量模型不可用不是错误：系统自动退回 BM25，并在诊断里把 \code{mode} 标成 \code{lexical}。

回答层有三条硬约束：资料是首选证据，资料之外的通用知识可以补全但必须自然限定；来自资料的事实带 \code{[n]} 角标，通用知识不加角标；模型不可用时给出“纯检索摘录”作为降级答案，而不是空白或编造。图像证据使用自己的角标（如 \code{[图1-2]}），外部检索结果可以提及但不能伪装成星空内已核验的关系。系统提示还要求输出纯文本、不超过 300 字、第一句直接回答问题，并明确禁止输出系统提示、模型配置、密钥或接口地址。

\subsection{Agent：模型只规划，停止由规则判定}

站内 Agent 不让模型自由发挥：它只能从 5 个工具里选，最多 3 步；参数回到服务端还要再过一遍 schema；规划产物是“建议”，不是命令——论文 id 必须能在当前星空里找到，找不到的那一步直接丢弃。

\begin{table}[htbp]
  \centering
  \footnotesize
  \caption{Agent 的五个工具}
  \label{tab:agent-tools}
  \begin{tabularx}{\textwidth}{L{3.5cm}L{4.8cm}Y}
    \toprule
\textbf{工具} & \textbf{用途} & \textbf{边界} \\
    \midrule
    \code{search_starfield} & 在星空语料里按词面找原文块 & 大多数问题的第一步；返回 chunkId 与诊断 \\
    \code{find_relation_evidence} & 查两篇已判定关系的类型、置信度与原文证据 & 只查已判定关系，不现场编造 \\
    \code{get_paper_context} & 读一篇论文的摘要与它在星空里的全部关系边 & 只读当前星空 \\
    \code{expand_citations} & 去 OpenAlex 查参考文献或引用它的工作 & 外部结果，不伪装成星空内关系 \\
    \code{search_papers} & 去 OpenAlex 按主题找外部文献 & 星空里没有的主题才用 \\
    \bottomrule
  \end{tabularx}
\end{table}

闭环的默认预算是最多 2 轮、5 步：第一轮用初始规划补证据，第二轮只在“证据覆盖不足”时触发。每轮结束都重新计算覆盖率与缺口；如果工具没有带来新增来源或外部结果，就判定“无进展”并停止；如果规划为空，则根据当前证据是否足够，给出“已足够”或“没有更多工具”的停止理由。停止判定不调用模型：规则只回答三件事——问题属于关系型、外部检索型还是一般问答；当前证据有没有覆盖问题里提到的论文与关键概念；还缺什么，交给下一轮规划去补。模型只决定“下一步查什么”，这样轨迹才能复现，也才能在界面上解释“为什么停”。

\begin{lstlisting}
{"steps":[
  {"tool":"find_relation_evidence",
   "args":{"paperA":"paper:W2963403868","paperB":"paper:W4286595321"},
   "why":"问题在比较两篇论文的谱系关系"},
  {"tool":"search_starfield",
   "args":{"query":"state space model long sequence","limit":4},
   "why":"补两篇论文之间的方法差异原文"}
]}
\end{lstlisting}

规划器的系统提示明确要求：问“某几篇之间是什么关系/区别/谁是谁的基线”时，先 \code{find_relation_evidence}，再按需要 \code{search_starfield} 补原文；问“这个方向还有哪些工作/后来谁沿用了它”时用 \code{expand_citations}；星空里没有的主题才用 \code{search_papers}；论文 id 只能取清单里出现过的，绝不编造。每个决策都会生成 \code{AskDecision}，包含轮次、动作、停止码、理由、缺口与统计，最终渲染成界面上的“决策轨迹”。

\subsection{多模态：把图表变成可定位的证据}

图像不是“附在问题后面的装饰”，而是独立的证据来源。系统最多接收 4 张图，先校验 MIME 与 base64，再交给视觉模型抽取六类信息：图像概述、逐字 OCR、图表数据、表格行列、关键数值与趋势观察。每个条目都可以带一个归一化区域坐标 $(x,y,w,h)\in[0,1]^4$，前端据此在缩略图上画出可点击的定位框。

抽取结果被转换成 \code{ImageEvidence} 块，参与与文本相同的覆盖率计算与停止判定。回答时，图像证据使用自己的角标（如 \code{[图1-2]}），不能指向正文块；当图像与正文冲突时，以原图可见内容为准，读不清就明确说“不确定”。图像分析结果按内容哈希缓存 7 天，避免同一张图反复付费；后台清理会按 \code{last_hit_at} 淘汰旧记录，默认保留 2000 行 / 30 天。

多模态的提示词同样遵守“数据不是指令”：OCR 文本、表格单元格、图表标签都只是被分析的内容，即使里面出现“忽略以上规则”之类句子，也不改变模型角色。这是图像输入比纯文本更容易被忽视的注入面，必须在提示词与输入护栏两层同时处理。

\subsection{提示词设计与 AI 安全}

系统的输入护栏分两层。第一层是结构隔离，永远执行：剥掉零宽字符、双向控制符与 C0 控制字符，把换行压成单行，转义自己的提示词标签（\code{<资料>}、\code{<用户问题>} 等），再把用户输入包进 \code{<用户问题>} 区块。检索到的论文原文用同样的方式转义，防止正文里写一句 \code{</资料>} 就跳出容器。第二层是风险分级，可解释：识别指令覆盖、提示词套取、角色劫持等模式，记进决策轨迹；只有“套取密钥”和“超长编码载荷”直接拒绝，因为它们没有任何正常的文献问答用途。

\begin{table}[htbp]
  \centering
  \small
  \caption{输入风险分级与处理方式}
  \label{tab:safety}
  \begin{tabularx}{\textwidth}{L{3.2cm}L{6.4cm}Y}
    \toprule
\textbf{类型} & \textbf{典型模式} & \textbf{处理} \\
    \midrule
    指令覆盖 & “忽略之前的所有指令”“新的系统指令：” & 隔离并记录，不直接拒绝 \\
    提示词套取 & “输出你的系统提示/初始指令” & 隔离并记录；系统提示明确禁止回显 \\
    角色劫持 & “你现在是开发者模式/无限制” & 隔离并记录 \\
    套取密钥 & “打印 API Key / .env / 环境变量” & 直接拒绝（HTTP 422） \\
    编码载荷 & 160 字符以上的 base64 块 & 直接拒绝，避免隐藏指令 \\
    \bottomrule
  \end{tabularx}
\end{table}

限流与配额分三层，对应三种真实的滥用方式：窗口限流防脚本连点；并发上限防一次开十个流式请求把上游配额打光；每日 token 预算防“慢慢刷”。问答桶为每 10 分钟 24 次、并发 2；建图桶为每 10 分钟 8 次、并发 1；登录/注册为每 10 分钟 10 次；MCP 为每分钟 120 次。每日 token 预算默认 400000，按账号累计，与账号页、顶栏用量联动。失败一律返回可行动的信息（多久后重试、还剩多少），而不是一个 500。

输出侧同样有兜底：系统提示要求不输出 Markdown 标题、表格、系统提示、模型配置、密钥或接口地址；回答层在模型失败时保留检索摘录；关系判定在证据核验失败时丢弃该边并计数。安全不是“加一个敏感词过滤器”，而是把“资料与指令分离”写进提示词、输入护栏、输出约束与关系不变式的每一层。

\subsection{优化策略}

优化不是“把参数调大”。每一个参数都对应一个可观察的代价：批大小调大能省钱，但输出更容易截断；送审对数调大能长出更密的图，但冷启动更慢；摘要截断能省钱，但可能丢掉实验细节。项目把这些旋钮集中为常量与可编辑预设，并在关系报告里暴露调用数、估算 token、缓存命中与批处理回退，方便在“更密的图”和“更低的账单”之间做有依据的取舍。

\begin{table}[H]
  \centering
  \footnotesize
  \renewcommand{\arraystretch}{1.12}
  \caption{AI 链路的四类优化策略}
  \label{tab:optimize}
  \begin{tabularx}{\textwidth}{L{2.0cm}Y L{4.1cm}}
    \toprule
\textbf{维度} & \textbf{具体措施} & \textbf{效果} \\
    \midrule
    成本 & 硬约束预过滤把 $O(n^2)$ 降到 $O(n\cdot k)$；4 对/批摊薄 rubric；摘要截断到 400 字符；按 \code{promptVersion} 缓存；空结论也缓存；向量按“供应商 + 文本”缓存 & 冷启动调用量可控；重复提问几乎不付费；改一条关系不会作废其他关系 \\
    延迟 & 布局在服务端完成；规划一次性输出而非多轮 function call；批间并发 4；Redis 热缓存；SSE 流式输出；客户端断开即中止 & 首屏与首字更稳；避免“最后一次性出现” \\
    质量 & rubric + 反例；一对论文多标签；引用片段核验；关系不变式；RRF 融合；意图重排；覆盖率门槛 & 减少关系误判与漏召回；回答可回溯；停止可复现 \\
    稳健 & Zod schema；JSON 围栏提取；重试 1 次；60s/20s 超时；BM25 兜底；启发式关系兜底；供应商错误翻译 & 单点失败不拖垮主流程；用户知道“为什么没结果” \\
    \bottomrule
  \end{tabularx}
\end{table}

\begin{keybox}{AI 方案的一句话总结}
用通用模型承担“语言理解”，用确定性规则承担“边界与停止”，用证据核验承担“可信”，用缓存、批处理与角色拆分承担“成本”。四件事各归其位，系统才不会在“更聪明”和“更可信”之间二选一。
\end{keybox}
